% Generated by tests/export_latex.pl from the knowledge base. Do not edit by hand.
\begin{footnotesize}
\begin{longtable}{L{2.8cm}L{1.3cm}L{8.2cm}L{2.0cm}}
  \caption{Maintenance-action mapping, generated from \texttt{kb/maintenance.pl}: the corrective plan issued for each diagnosis, in order.}\label{tab:maintenance}\\
  \toprule
  Diagnosis & Priority & Ordered actions & Source \\
  \midrule
  \endfirsthead
  \toprule
  Diagnosis & Priority & Ordered actions & Source \\
  \midrule
  \endhead
  \texttt{crah\_\allowbreak{}fan\_\allowbreak{}failure} & urgent & 1.~Check the fan drive fault log and the motor current \newline 2.~Inspect the motor, belt or coupling, and bearings; repair or replace the failed part \newline 3.~Bring redundant cooling online or raise the output of neighbouring units during the repair \newline 4.~After repair, confirm fan power matches the affinity-law prediction and rack inlets fall below 27 C & \texttt{ashrae\_\allowbreak{}hb\_\allowbreak{}fans} \\
  \texttt{chw\_\allowbreak{}supply\_\allowbreak{}loss} & urgent & 1.~Confirm chilled-water plant status: pumps running and isolation valves open \newline 2.~Compare commanded and actual valve stroke; repair or replace a seized actuator \newline 3.~Vent trapped air from the coil \newline 4.~Confirm water flow matches the valve position and supply air is back at setpoint & \texttt{a4\_\allowbreak{}chw\_\allowbreak{}flow} \\
  \texttt{hot\_\allowbreak{}air\_\allowbreak{}recirculation} & high & 1.~Inspect the rack for missing blanking panels and unsealed cable openings \newline 2.~Check containment doors and roof panels for gaps \newline 3.~Install blanking panels and seal the openings \newline 4.~Re-measure top and middle inlet temperatures to confirm the difference is below 3 K & \texttt{eu\_\allowbreak{}coc\_\allowbreak{}dc} \\
  \texttt{normal} & routine & 1.~No corrective action. Continue routine monitoring. & \texttt{ashrae\_\allowbreak{}tc99\_\allowbreak{}2021} \\
  \texttt{unexplained\_\allowbreak{}abnormality} & high & 1.~Verify the affected sensors before acting \newline 2.~Escalate to a technician for manual investigation & \texttt{closed\_\allowbreak{}world} \\
  \texttt{filter\_\allowbreak{}clogging} & high & 1.~Replace the CRAH air filters \newline 2.~Check the filter differential pressure gauge against a reference \newline 3.~Review the filter replacement interval and recent dust sources such as building work \newline 4.~After replacement, confirm filter pressure drop and fan speed are back to normal & \texttt{a8\_\allowbreak{}filter} \\
  \texttt{server\_\allowbreak{}fan\_\allowbreak{}failure} & high & 1.~Check the server management controller event log for the failed fan module \newline 2.~Migrate or reduce the workload on the server until it is repaired \newline 3.~Hot-swap the failed fan module \newline 4.~Confirm CPU temperature and clock frequency return to normal & \texttt{a9\_\allowbreak{}server\_\allowbreak{}fan} \\
  \texttt{branch\_\allowbreak{}overload} & urgent & 1.~Identify the servers on the overloaded phase from the rack power records \newline 2.~Rebalance server power supplies across the three phases \newline 3.~Check for equipment added without a capacity review \newline 4.~Confirm every phase is below 80\% of its breaker rating & \texttt{nec\_\allowbreak{}continuous} \\
  \texttt{battery\_\allowbreak{}degradation} & high & 1.~Confirm with a capacity discharge test as IEEE Std 1188 recommends \newline 2.~Check battery room temperature and its cooling \newline 3.~Replace weak blocks or the whole string, depending on the test result \newline 4.~Record a new resistance baseline after replacement & \texttt{a11\_\allowbreak{}battery} \\
  \texttt{ups\_\allowbreak{}on\_\allowbreak{}bypass} & urgent & 1.~Read the UPS alarm log to find why it transferred to bypass \newline 2.~Call the UPS service engineer to repair the inverter or rectifier fault \newline 3.~Postpone any upstream electrical work while the load has no battery protection \newline 4.~Return the UPS to double conversion after the repair and confirm normal efficiency & \texttt{a12\_\allowbreak{}ups} \\
  \texttt{sensor\_\allowbreak{}bias} & high & 1.~Check the sensor against a calibrated reference instrument \newline 2.~Check the sensor mounting and position \newline 3.~Recalibrate or replace the sensor \newline 4.~Do not act on readings from this sensor until it is verified & \texttt{a13\_\allowbreak{}neighbour} \\
  \texttt{telemetry\_\allowbreak{}dropout} & high & 1.~Check the monitoring agent, controller and network link for this telemetry group \newline 2.~Restore the data feed \newline 3.~Treat the equipment status as unknown until fresh data arrives \newline 4.~Inspect critical equipment on site if the feed cannot be restored quickly & \texttt{a14\_\allowbreak{}telemetry} \\
  \bottomrule
\end{longtable}
\end{footnotesize}
